\inicializarSeccion{Dielectroforesis}

\tab La dielectroforesis (DEP) es una técnica que permite mover y manipular las partículas diminutas, como lo son las células de nuestra sangre, usando los campos eléctricos. A diferencia de otros métodos que necesitan de partículas con carga eléctrica, la DEP funciona incluso con células neutras \parencite{Dielectrophoresis_2026}. Lo increíble de esta técnica es que no necesita de químicos o colorantes costosos para identificar las enfermedades; reconoce las células por su estructura física y estado de salud natural \parencite{Dielectrophoresis_2026}.

\tab Este proceso funciona creando un campo eléctrico no uniforme, un espacio donde la fuerza eléctrica es más intensa en unos puntos que en otros. Cuando el glóbulo rojo está en el campo eléctrico, las cargas internas se reordenan, creando un dipolo que interactúa con el campo externo \parencite{GiduThuri}.

\tab Las fuerzas que se sienten en cada uno de los glóbulos rojos no son iguales para todos, ya que dependen de factores como el tamaño de la célula, la intensidad del campo y sus propiedades dieléctricas \parencite{Dielectrophoresis_2026}. Estas propiedades determinan la respuesta de la célula: dependiendo de la frecuencia eléctrica, esta puede ser atraída a zonas con mayor o menor intensidad eléctrica \parencite{GiduThuri}.

\tab En el diagnóstico de la malaria, esta técnica es especialmente útil porque el parásito altera la firma eléctrica del glóbulo rojo afectado \parencite{Dielectrophoresis_2026}. Debido a estos cambios en la estructura celular, los glóbulos rojos sanos e infectados se comportan de manera distinta bajo el mismo campo eléctrico, lo que permite separar las células enfermas de las sanas y ayudar en la detección de la enfermedad \parencite{Dielectrophoresis_2026}.
